%%%%%%%%%%%%%%%%%%%%%%%%%%%%%%%%%%%%%%%%%%%%%%%%%%%%%%%%%%%%%%%%%%%%%%
% Lớp 10 | Phương pháp tọa độ trong mặt phẳng (𝑂𝑥𝑦) | Đường tròn trong mặt phẳng tọa độ
% Dạng 0H9?4-1: Tìm tâm, bán kính và điều kiện là đường tròn
% Số câu: 31  (sinh tự động bởi tach_cau.py; từ nay thầy sửa trực tiếp file này)
%%%%%%%%%%%%%%%%%%%%%%%%%%%%%%%%%%%%%%%%%%%%%%%%%%%%%%%%%%%%%%%%%%%%%%

%%%%%%%%%%%%%%%% TRẮC NGHIỆM 4 PHƯƠNG ÁN %%%%%%%%%%%%%%%%

% nguồn: MaHoaNH_TracNghiem_Lop10
\begin{ex}%[0H9N4-1]%[ma:be8ad567]
	Trong mặt phẳng tọa độ $Oxy$, cho đường tròn $(C)\colon (x-1)^2+(y-2)^2=9$. Tìm tọa độ  tâm $I$ và bán kính $R$ của đường tròn $(C)$.
	\choice 
		 {$I(1; 2)$ và $R=9$}
		{\True $I(1; 2)$ và  $R=3$}
		 {$I(-1;-2)$ và $R=9$}
		{$I(-1;-2)$  và $R=3$}
	\loigiai{ 
	Đường  tròn $(C)$ có tâm $I(1;2)$ và bán kính $R=3$.}
\end{ex}

% nguồn: MaHoaNH_TracNghiem_Lop10
\begin{ex}%[0H9N4-1]%[ma:1c45695b]
	Đường tròn $2x^2+2y^2-8x+4y-1=0$ có tâm là  điểm nào sau đây?
	\choice 
		 {$(-8; 4)$}
		{\True  $(2;-1)$}
		{$(8;-4)$} 
		{$(-2;  1)$}
	\loigiai{$2x^2+2y^2-8x+4y-1=0\Leftrightarrow x^2+y^2-4x+2y-\dfrac{1}{2}=0\Leftrightarrow  (x-2)^2+(y+1)^2=\dfrac{11}{2}$.\\
		Vậy tâm của đường tròn đã  cho là $I(2;-1)$.
	} 
\end{ex}

% nguồn: MaHoaNH_TracNghiem_Lop10
\begin{ex}%[0H9N4-1]%[ma:b1f0acb3]
	Trong mặt phẳng toạ độ $Oxy$, phương trình đường tròn có tâm  $I(-3;2)$ và bán kính $R=3$ là
	\choice 
		{$(x+3)^2+(y-2)^2=3$} 
		 {$(x-3)^2+(y+2)^2=9$}
		 {$(x-3)^2+(y+2)^2=3$}
		 {\True $(x+3)^2+(y-2)^2=9$}
	\loigiai{ 
		Phương trình đường tròn có tâm $I(-3;2)$ và bán kính  $R=3$ là $$[x-(-3)]^2+(y-2)^2=3^2 \Leftrightarrow (x+3)^2+(y-2)^2=9.$$
	} 
\end{ex}

% nguồn: MaHoaNH_TracNghiem_Lop10
\begin{ex}%[0H9N4-1]%[ma:f791ae7c]
	Phương trình nào sau đây là phương trình đường tròn tâm  $I(a;b)$ và bán kính $R$?
	\choice 
		{$( x+a)^2-(y+b)^2=R^2$} 
		{\True $(  x-a)^2+(y-b)^2=R^2$}
		{$x^2+y^2-2ax+2by+(a^2+b^2-R^2)=0$} 
		{$x^2-y^2-2ax-2by+(a^2+b^2-R^2)=0$} 
	\loigiai{ 
		Phương trình đường tròn tâm $I(a;b)$, bán kính $R$ là $$(x-a)^2+(y-b)^2=R^2.$$ 
	} 
\end{ex}

% nguồn: MaHoaNH_TracNghiem_Lop10
\begin{ex}%[0H9N4-1]%[ma:fc6487ba]
	Tìm tọa độ tâm $I$ của  đường tròn $(C)\colon(x-2)^2+(y+5)^2=9$.
	\choice 
		{\True  $I(2;-5)$}
		 {$I(2;5)$}
		 {$I(-2;5)$}
		{$I(-2;-5)$} 
	 \loigiai{
		Tọa độ tâm của  đường tròn $(C)$ là $I(2;-5)$.
	 }
\end{ex}

% nguồn: MaHoaNH_TracNghiem_Lop10
\begin{ex}%[0H9N4-1]%[ma:cfa4353b]
	Trong  mặt phẳng $Oxy$, đường tròn $(C)\colon x^2+y^2-2 x+10 y+17=0$ có tọa độ tâm $I$ là
	 \choice
		 {$I(-1;-5)$}
		{\True $I(1;-5)$} 
		 {$I(1;5)$}
		 {$I(-1;5)$}
	\loigiai{ 
		Đường tròn $(C)$ có  tâm là $I(1;-5)$.
	 }
\end{ex}

% nguồn: MaHoaNH_TracNghiem_Lop10
\begin{ex}%[0H9N4-1]%[ma:06364874]
	Cho  đường tròn $\left(C\right)\colon (x-2)^2+(y+6)^2=8$ có tâm $I$ và bán kính là
	 \choice
		{\True  $I(2;-6),R=2\sqrt{2}$}
		 {$I(2;-6),R=2$}
		{$I(2;6),R=2\sqrt{2}$} 
		 {$I(2;-6),R=8$}
	\loigiai{ 
		$(C)$ có tâm $I(2;-6)$ và bán kính $R=\sqrt{8}=2\sqrt{2}$ 
	 }
\end{ex}

% nguồn: MaHoaNH_TracNghiem_Lop10
\begin{ex}%[0H9N4-1]%[ma:32ff2fbd]
	Trong mặt phẳng $Oxy$, cho đường tròn $(C)\colon \left(x-1\right)^2+\left(y+3\right)^2=16$. Tọa độ tâm $I$ của đường  tròn $(C)$ là
	 \choice
		{\True  $\left(1;-3\right)$}
		{$\left(-3;1\right)$} 
		 {$\left(-1;-3\right)$}
		 {$\left(1;3\right)$}
	\loigiai{ 
	Tọa độ tâm $I$ của  đường tròn $(C)$ là $(1; -3)$.
	} 
\end{ex}

% nguồn: MaHoaNH_TracNghiem_Lop10
\begin{ex}%[0H9N4-1]%[ma:234bcfc1]
	Trong mặt phẳng $Oxy$, xác định tâm và bán kính của đường  tròn
	\par $(C)\colon (x-2)^{2}+(y+5)^{2}=6$. 
	\choice 
		{\True $I(2;-5)$, $R=\sqrt{6}$} 
		{$I(-2;5)$, $R=\sqrt{6}$} 
		{$I(5;-2)$,  $R=\sqrt{6}$}
		{$I(-5;2)$, $R=\sqrt{6}$} 
	 \loigiai{
		Tâm $I$ và bán kính $R$ của đường tròn $(C)\colon (x-2)^{2}+(y+5)^{2}=6$  là $I(2;-5)$ và bán kính $R=\sqrt{6}$.
	} 
\end{ex}

% nguồn: MaHoaNH_TracNghiem_Lop10
\begin{ex}%[0H9N4-1]%[ma:b86a4974]
	Một đường tròn có  tâm $I(3;-2)$ tiếp xúc với đường thẳng $\Delta\colon x-5y+1=0$. Bán kính đường tròn bằng
	\choice 
		{\True $\dfrac{14}{\sqrt{26}}$} 
		{$\dfrac{7}{13}$} 
		{$\sqrt{26}$} 
		 {$6$}
	 \loigiai{
		 Gọi bán kính của đường tròn là $R$.\\
		Khi đó $R=\mathrm{d}(I,  \Delta)=\dfrac{|3-5\cdot(-2)+1|}{\sqrt{1^2+(-5)^2}}=\dfrac{14}{\sqrt{26}}$.
	 }
\end{ex}

% nguồn: MaHoaNH_TracNghiem_Lop10
\begin{ex}%[0H9N4-1]%[ma:e6c59ef0]
	Cho đường cong $(C)\colon (x-1)^2+y^2-m-5=0$. Tìm  $m$ để $(C)$ là một phương trình đường tròn
	\choice 
		{\True $m>-5$} 
		 {$m \geq -5$}
		 {$m<-5$}
		{$m \leq  -5$}
	\loigiai{ 
		 Ta có $(x-1)^2+y^2-m-5=0 \Leftrightarrow (x-1)^2+y^2=m+5$.\\
		$(C)$ là phương trình đường tròn $\Leftrightarrow m+5 > 0 \Leftrightarrow m>-5$. 
	 }
\end{ex}

% nguồn: MaHoaNH_TracNghiem_Lop10
\begin{ex}%[0H9N4-1]%[ma:ab0935fe]
	Phương trình nào dưới  đây là phương trình đường tròn?
	\choice 
		{$x^2+2y^2-4x-8y+1=0$} 
		{$4x^2+y^2-10x-6y-2=0$} 
		 {$x^2+y^2-2x-8y+20=0$}
		{\True  $x^2+y^2-4x+6y-12=0$}
	\loigiai{ 
		Xét phương trình $x^2+y^2-4x+6y-12=0$  có $a=2$; $b=-3$; $c=-12$;\\
		 suy ra $a^2+b^2-c=4+9+12=25>0$ nên $x^2+y^2-4x+6y-12=0$ là phương trình đường tròn.
	 }
\end{ex}

% nguồn: MaHoaNH_TracNghiem_Lop10
\begin{ex}%[0H9N4-1]%[ma:b3d9503c]
	Trong mặt phẳng toạ độ $Oxy$, cho đường tròn  $(C) \colon x^2+y^2-6x+4y+1=0$. Tâm $I$ của $(C)$ có tọa độ là
	\choice 
		{\True  $I(3;-2)$}
		{$I(6;-4)$} 
		{$I(-6;4)$} 
		 {$I(-3;2)$}
	\loigiai{ 
		Ta có  $x^2+y^2-6x+4y+1=0 \Leftrightarrow x^2+y^2-2\cdot 3x-2\cdot (-2)y+1=0$.\\
		Vậy tâm  $I(3;-2)$.
	} 
\end{ex}

% nguồn: MaHoaNH_TracNghiem_Lop10
\begin{ex}%[0H9N4-1]%[ma:1c40d516]
	Trong mặt phẳng $Oxy$, phương trình nào  sau đây là phương trình của đường tròn?
	\choice 
		 {$x^2+2y^2-4x-8y+1=0$}
		{\True  $x^2+y^2-4x+6y-12=0$}
		 {$x^2+y^2-2x-8y+20=0$}
		 {$4x^2+y^2-10x-6y-2=0$}
	\loigiai{ 
		Để là  phương trình đường tròn thì điều kiện cần là hệ số của $x^2$ và $y^2$ phải bằng nhau nên loại được đáp án $\mathrm{\textbf{A}}$ và $\mathrm{\textbf{D}}$.\\
		Ta có $x^2+y^2-2x-8y+20=0 \Leftrightarrow (x-1)^2+(y-4)^2+3=0$  vô lý.\\
		Ta  có $x^2+y^2-4x+6y-12=0 \Leftrightarrow (x-2)^2+(y+3)^2=25$ là phương trình đường tròn tâm $I(2;-3)$, bán kính $R=5$.
	} 
\end{ex}

% nguồn: MaHoaNH_TracNghiem_Lop10
\begin{ex}%[0H9N4-1]%[ma:2d9e058b]
	Tâm  của đường tròn $(C)$ có phương trình $(x-3)^{2}+(y+4)^{2}=12$ là
	 \choice
		 {\True $I(3;-4)$}
		{$I(3;4)$} 
		 {$I(-3;4)$}
		 {$I(4;3)$}
	\loigiai{ 
		Phương  trình đường tròn $(x-a)^{2}+(y-b)^{2}=R^{2}$ với $I(a,b)$ là tâm của đường tròn.\\
		Từ  phương trình tổng quát của $(C)\colon (x-3)^{2}+(y+4)^{2}=12$ ta suy ra $a=3$, $b=-4$.\\
		Vậy tâm của đường tròn $(C)$  là $I(3;-4)$.
	} 
\end{ex}

% nguồn: MaHoaNH_TracNghiem_Lop10
\begin{ex}%[0H9H4-1]%[ma:21749615]
	Trong mặt phẳng $Oxy$, cho đường cong $\left(C_m \right) \colon x^2+y^2-8x+10y+m=0$. Với giá trị nào của $m$  thì $\left(C_m \right)$ là đường tròn có bán kính bằng $7$?
	 \choice
		 {$m=4$}
		{$m=8$} 
		 {\True $m=-8$}
		 {$m =-4$}
	 \loigiai{
		Ta có $R=\sqrt{4^2+5^2-m}=7\Leftrightarrow m=-8$. 
	 }
\end{ex}

% nguồn: MaHoaNH_TracNghiem_Lop10
\begin{ex}%[0H9H4-1]%[ma:57122048]
	Trong mặt phẳng $Oxy$, đường tròn $(C)$ tâm $I(-4;3)$ và tiếp xúc  với trục tung có phương trình là
	\choice 
		 {$x^2+y^2-4x+3y+9=0$}
		{\True $(x+4)^2+(y-3)^2=16$} 
		 {$(x-4)^2+(y+3)^2=16$}
		 {$x^2+y^2+8x-6y-12=0$}
	\loigiai{ 
		Vì $(C)$  tiếp xúc với $y'Oy$ và có tâm $I(-4;3)$ nên $a=-4$, $b=3$, $ R=\left|a\right|=4$.\\
		Do đó, $(C)$ có phương trình $\left(x+4\right)^2+(y-3)^2=16$. 
	} 
\end{ex}

% nguồn: MaHoaNH_TracNghiem_Lop10
\begin{ex}%[0H9H4-1]%[ma:d849913a]
	Trong các phương  trình sau, phương trình nào là phương trình của một đường tròn?
	 \choice
		{$5x^2+4y^2-x+4y-1=0$} 
		 {$x^2+y^2+2x-4y+9=0$}
		 {\True $x^2+y^2-4x-2y+1=0$}
		{$x^2+y^2-6x+4y+13=0$} 
	\loigiai{ 
		Xét  phương trình $x^2+y^2-4x-2y+1=0$
		ta có $a^2 + b^2 - c = 2^2 + 1^2  -1 = 4 > 0$ nên $x^2+y^2-4x-2y+1=0$ là phương trình của một đường tròn.
	 }
\end{ex}

% nguồn: MaHoaNH_TracNghiem_Lop10
\begin{ex}%[0H9H4-1]%[ma:cf2f2c92]
	Trong mặt phẳng $Oxy$, tâm của đường tròn $2x^2+2y^2-8x+4y-1=0$ có  tọa độ là
	\choice 
		{$(-8;4)$} 
		{\True  $(2;-1)$}
		 {$(8;-4)$}
		 {$(-2;1)$}
	\loigiai{ 
		Ta có 
		\[ 
		2x^2+2y^2-8x+4y-1=0 \Leftrightarrow  x^2+y^2-4x+2y-\dfrac{1}{2}=0.
		 \]
		Vậy tâm  là $I(2;-1)$.
	 }
\end{ex}

% nguồn: MaHoaNH_TracNghiem_Lop10
\begin{ex}%[0H9H4-1]%[ma:31f86204]
	Đường tròn  $3x^2+3y^2-6x+9y-9=0$ có bán kính bằng
	 \choice
		 {$\dfrac{15}{2}$}
		{\True $\dfrac{5}{2}$} 
		 {$25$}
		 {$\sqrt{5}$}
	\loigiai{Ta  có
		 \[
		 3x^2+3y^2-6x+9y-9=0\Leftrightarrow x^2+y^2-2x+3y-3=0.
		\] 
		 Suy ra $P=1^2+\left(\dfrac{3}{2} \right)^2-\left(-3\right)=\dfrac{25}{4} > 0$.\\
		Vậy bán kính  là $R=\dfrac{5}{2}$.
	} 
\end{ex}

% nguồn: MaHoaNH_TracNghiem_Lop10
\begin{ex}%[0H9H4-1]%[ma:ee17d75a]
	Trong mặt phẳng $Oxy$, đường tròn đi qua $3$ điểm $A(11; 8)$; $ B(13;8)$; $C(14;7)$ có bán kính  $R$ bằng
	\choice 
		{$2$} 
		{$1$} 
		 {\True $\sqrt{5}$}
		 {$\sqrt{2}$}
	\loigiai{ 
		Gọi phương trình đường tròn cần tìm có dạng $x^2+y^2-2ax-2by+c=0 \left(a^2+b^2-c > 0\right)$.\\ 
		Đường tròn đi qua $3$ điểm $A(11; 8)$, $B(13;8)$, $C(14;7)$  nên ta có\\
		$\heva{&121+64-22a-16b+c=0\\&169+64-26a-16b+c=0 \\ &196+49-28a-14b+c=0} \Leftrightarrow  \heva{&a=12\\&b=6\\&c=175.}$\\
		Ta có  $R=\sqrt{a^2+b^2-c}=\sqrt{5}$.\\
		Vậy phương trình đường tròn đi qua $3$ điểm $A(11; 8)$, $B(13;8)$, $C(14;7)$ có  bán kính là $R=\sqrt{5}$.
	} 
\end{ex}

% nguồn: MaHoaNH_TracNghiem_Lop10
\begin{ex}%[0H9Y4-1]%[ma:8efc15da]
	Cho đường tròn $(C)\colon (x-1)^2+(y-2)^2=25$.  Đường tròn $(C)$ có
	\choice 
		{Tâm $I(1;2)$ và bán kính  $R=25$}
		{Tâm $I(-1;-2)$ và bán kính  $R=25$}
		 {\True Tâm $I(1;2)$ và bán kính $R=5$}
		{Tâm  $I(-1;-2)$ và bán kính $R=5$}
	 \loigiai{
		 Đường tròn $(C)$ có tâm $I(1;2)$ và bán kính $R=\sqrt{25}=5$.
	} 
\end{ex}

% nguồn: MaHoaNH_TracNghiem_Lop10
\begin{ex}%[0H9Y4-1]%[ma:59e250f8]
	Cho đường tròn $(C)\colon x^2+y^2+6x-4y+2=0$. Đường tròn  $(C)$ có
	 \choice
		{Tâm $I(-3;2)$ và bán kính $R=11$} 
		{\True Tâm  $I(-3;2)$ và bán kính $R=\sqrt{11}$}
		 {Tâm $I(3;-2)$ và bán kính $R=11$}
		{Tâm $I(3;-2)$ và  bán kính $R=\sqrt{11}$}
	\loigiai{ 
		Ta có  $a=\dfrac{6}{-2}=-3$ và $b=\dfrac{-4}{-2}=2$ và $c=2$.\\
		 $\Rightarrow a^2+b^2-c=(-3)^2+2^2-2=11$.\\
		$\Rightarrow$ Đường tròn $(C)$ có tâm  $I(-3;2)$ và bán kính $R=\sqrt{a^2+b^2-c}=\sqrt{11}$.
	 }
\end{ex}


%%%%%%%%%%%%%%%% ĐÚNG/SAI %%%%%%%%%%%%%%%%

% nguồn: MaHoaNH_DungSai_Lop10
\begin{ex}%[0H9N4-1]%[ma:2717d8b7]
	Các khẳng định sau  đây đúng hay sai?
	\choiceTF 
		{Đường tròn $(C)\colon (x-1)^2+(y+2)^2=4$ có tâm $I(1;2)$ và bán kính $R=2$} 
		{Phương trình $x^2+y^2+2=0$ là phương trình của  một đường tròn}
		{\True Đường tròn tâm  $O(0;0)$ và đi qua điểm $M(1;2)$ có bán kính $R=\sqrt{5}$}
		{\True Đường tròn đường kính  $AB$ với $A(-2;-1)$, $B(4;3)$ có tâm $J(1;1)$}
	 \loigiai{
		\begin{itemchoice}
			\itemch Đường tròn $(C)\colon (x-1)^2+(y+2)^2=4$ có tâm  $I(1;-2)$ và bán kính $R=2$.
			\itemch Phương trình $x^2+y^2+2=0$ không phải là phương trình của một đường tròn  vì $x^2+y^2=-2<0$.
			\itemch Đường tròn tâm $O(0;0)$ và đi qua điểm $M(1;2)$  có bán kính $R=OM=\sqrt{1^2+2^2}=\sqrt{5}$.
			\itemch Đường tròn đường kính $AB$  với $A(-2;-1)$, $B(4;3)$ có tâm $J(1;1)$ là trung điểm của $AB$.
		\end{itemchoice}
	 }
\end{ex}

% nguồn: MaHoaNH_DungSai_Lop10
\begin{ex}%[0H9N4-1]%[ma:22af6f93]
	Trong mặt phẳng tọa độ $Oxy$, cho đường tròn $(C)\colon (x+1)^2+y^2=1$, điểm $A(2;1)$ và $B(1;3)$. Các  khẳng định sau đây đúng hay sai?
	 \choiceTF
		{Đường tròn $(C)$ có tâm  $I(1;0)$}
		{\True  Đường tròn $(C)$ có bán kính $R=1$}
		{\True Đường tròn tâm $A$ và đi  qua điểm $B$ có bán kính $R=\sqrt{5}$}
		{Đường tròn đường kính $AB$ có bán kính  $R=\sqrt{5}$}
	 \loigiai{
		\begin{itemize}
			\item Đường tròn $(C)$  có tâm $I(-1;0)$.
			\item Đường tròn $(C)$ có bán  kính $R=1$.
			\item Đường  tròn tâm $A$ và đi qua điểm $B$ có bán kính $R=AB=\sqrt{(1-2)^2+(3-2)^2}=\sqrt{5}$.
			\item Đường tròn  đường kính $AB$ có bán kính $R=\dfrac{AB}{2}=\dfrac{\sqrt{5}}{2}$.
		\end{itemize}
	} 
\end{ex}

% nguồn: MaHoaNH_DungSai_Lop10
\begin{ex}%[0H9H4-1]%[ma:583399a7]
	Trong các khẳng định sau đây, khẳng định nào đúng,  khẳng định nào sai?
	\choiceTF 
		{\True Phương trình $x^2+y^2+2x-6y+6=0$ là phương  trình của một đường tròn với bán kính $R=2$}
		{\True Phương trình $x^2+y^2-2(m+1)x+4y+2m^2+2=0$  là phương trình của một đường tròn khi và chỉ khi $-1<m<3$}
		{Đường tròn tâm $A(1;2)$, tiếp  xúc với đường thẳng $d\colon 2x+y-1=0$ có bán kính $R=3$}
		{Đường tròn có tâm  $I$ nằm trên trục $Ox$ và đi qua hai điểm $M(1;2)$, $N(4;-1)$ có bán kính $R=3$}
	\loigiai{ 
		\begin{itemchoice}
			\itemch Ta có $x^2+y^2+2x-6y+6=0 \Leftrightarrow (x+1)^2+(y-3)^2=4$  là phương trình của một đường tròn với bán kính $R=2$.
			\itemch Phương trình $x^2+y^2-2(m+1)x+4y+2m^2+2=0$ là phương trình của một đường tròn khi và chỉ khi $$(m+1)^2+(-2)^2-2m^2-2>0 \Leftrightarrow -m^2+2m+3>0 \Leftrightarrow -1<m<3.$$ 
			\itemch Đường tròn tâm $A(1;2)$, tiếp xúc với đường thẳng $d\colon 2x+y-1=0$ có bán kính  $$R=\mathrm{d}(A,d) = \dfrac{|2\cdot 1+2-1|}{\sqrt{2^2+1^2}} = \dfrac{3\sqrt{5}}{5}.$$
			\itemch Đường tròn có tâm $I$ nằm trên trục $Ox$ nên $I(m;0)$.\\ 
			Đường đi qua hai điểm $M(1;2)$, $N(4;-1)$  nên
			 $$MI=NI \Leftrightarrow \sqrt{(m-1)^2+4}=\sqrt{(m-4)^2+1} \Leftrightarrow m^2-2m+5=m^2-8m+17 \Leftrightarrow m=2.$$
			Do đó bán kính đường tròn  là $R=MI=\sqrt{(2-1)^2+4}=\sqrt{5}$.
		\end{itemchoice}
	 }
\end{ex}


%%%%%%%%%%%%%%%% TRẢ LỜI NGẮN %%%%%%%%%%%%%%%%

% nguồn: MaHoaNH_TraLoiNgan_Lop10
\begin{bt}%[0H9H4-1]%[ma:3837fdb7]
	Trong mặt phẳng $Oxy$, phương trình đường tròn tâm $I(5;6)$ và tiếp xúc với đường  thẳng $d\colon 3x-4y-6=0$ có bán kính bằng bao nhiêu?
	 \shortans{3}
	 \loigiai{
		Ta  có $R=\mathrm{d}(I;d)=\dfrac{|3.5-4.6-6|}{\sqrt{3^2+(-4)^2}}=3$.
	} 
\end{bt}

% nguồn: MaHoaNH_TraLoiNgan_Lop10
\begin{bt}%[0H9H4-1]%[ma:55950e4c]
	Biết rằng phương trình $x^2+y^2-2(m+2)x+4my+19m-6=0$ là một phương  trình đường tròn khi và chỉ khi $m\in (-\infty;a)\cup (b;+\infty)$, với $a < b$. Giá trị của $a + b$ bằng bao nhiêu?
	 \shortans{3}
	 \loigiai{
		Phương trình đã cho  là phương trình đường tròn khi và chỉ khi $$(m+2)^2+(-2m)^2-(19m-6) > 0\Leftrightarrow 5m^2-15m+10> 0\Leftrightarrow m\in (-\infty;1)\cup (2;+\infty).$$
		 Suy ra $a = 1$, $b = 2$. Vậy $a + b = 3$.
	 }
\end{bt}


%%%%%%%%%%%%%%%% TỰ LUẬN %%%%%%%%%%%%%%%%

% nguồn: MaHoaNH_TuLuan_Lop10
\begin{bt}%[0H9B4-1]%[ma:f760cccc]
	Trong mặt phẳng tọa độ $ Oxy $, cho elip có phương  trình $ \dfrac{x^2}{169}+\dfrac{y^2}{25} =1$. Tìm tọa độ các tiêu điểm, độ dài trục lớn và độ dài trục nhỏ của elip.
	 \loigiai{
		Ta  có $ \dfrac{x^2}{169}+\dfrac{y^2}{25} =1 \Rightarrow a=13, b=5 \Rightarrow c=\sqrt{169 -25}=12$.\\
		Suy ra $ (E) $ có các tiêu điểm $ F_1(-12;0) $ và $ F_2(12;0)  $.\\
		Độ dài  trục lớn $ 2a=26 $ và độ dài trục nhỏ $ 2b=10 $.
	 }
\end{bt}

% nguồn: MaHoaNH_TuLuan_Lop10
\begin{bt}%[0H9B4-1]%[ma:65caa0f3]
	Trong mặt phẳng tọa độ $Oxy$, lập phương trình chính tắc của Elip có độ dài trục lớn  là $8$, tiêu cự bằng $2\sqrt{5}$.
	\loigiai{ 
		Theo giả thiết, ta có $\heva{&2a = 8 \Rightarrow a = 4\\&2c = 2\sqrt{5}\Rightarrow c =  \sqrt{5}.}$\\
		Khi đó $b=\sqrt{a^2 - c^2}=\sqrt{16 - 5}=\sqrt{11}$.\\ 
		 Vậy phương trình Elip cần tìm là
		\[\dfrac{x^2}{a^2} + \dfrac{y^2}{b^2} = 1 \Leftrightarrow \dfrac{x^2}{16}+ \dfrac{y^2}{11}=1.\] 
	 }
\end{bt}

% nguồn: MaHoaNH_TuLuan_Lop10
\begin{bt}%[0H9B4-1]%[ma:74f37e2c]
	Trong mặt phẳng tọa  độ $Oxy$, cho Elip $(E)$ có phương trình $\dfrac{x^2}{9}+16y^2=1$. Tìm độ dài trục lớn, trục bé và các tiêu điểm của Elip.
	 \loigiai{
	\immini 
		 {
		Chuyển Elip  $(E)$ về dạng phương trình chính tắc
		$$\dfrac{x^2}{9}+16y^2=1\Leftrightarrow \dfrac{x^2}{9}+\dfrac{y^2}{\tfrac{1}{16}}=1.$$ 
		 }
		{ 
		\begin{tikzpicture}[scale=1.25,font=\footnotesize,line join=round,line cap=round,>=stealth]
				\draw (0,0) ellipse (3cm and  0.35cm);
				\draw[->] (-4,0) -- (4,0) node[below] {$x$}; 
				\draw[->]  (0,-1) -- (0,1) node[right] {$y$};
				\fill (0,0) circle  (1pt) node[below left]{$O$};
				\fill (2.8,0)  circle (1pt) node[below right]{$F_2$};
				\fill (-2.8,0) circle (1pt) node[below left]{$F_1$}; 
			\end{tikzpicture}
		 }\noindent
		Ta được  $a=3$, $b=\dfrac{1}{4}$ và $c=\sqrt{a^2-b^2}=\dfrac{\sqrt{143}}{4}$.\\
		 Khi đó
		\begin{itemize}
			\item Độ dài trục lớn là  $2a=6$.
			\item Độ dài trục bé  là $2b=\dfrac{1}{2}$.
			\item Các  tiêu điểm là $F_1\left( -\dfrac{\sqrt{143}}{4};0\right) $ và $F_2\left( \dfrac{\sqrt{143}}{4};0\right) $.
		\end{itemize}
	 }
\end{bt}
